\section{Discussion}
\label{sec:discussion}

This section interprets the scattering, scintillation, turbulence, and
energetics results sightline by sightline before returning to the
population-level claims; it is developed alongside those results.

\subsection{Physical interpretation}
\label{sec:disc-interpretation}

% TODO(disc-interpretation): Open with the small set of result statements that
% survive re-validation. State what each result means physically before making
% any sample-level claim. This paragraph should depend on, in order:
% 1. Section~\ref{sec:results-budget}: revalidated DM and scattering budgets.
% 2. Section~\ref{sec:dominant-systems}: trusted dominant foreground systems.
% 3. Section~\ref{sec:results-scintillation}: trusted scintillation constraints.
% 4. Section~\ref{sec:results-alpha}: trusted turbulence or morphology results.
% 5. Section~\ref{sec:burst-energies}: trusted energy measurements.
%
% Avoid restating the old revoked campaign results here. In particular, do not
% quote the retired host-dominated 10/11 comparison, the FRB 20230913A
% intervening attribution, scintillation-excess factors, rail-class tallies,
% alpha=4 limits, beta values, burst energies, or DM-budget residuals until the
% producing analysis has passed its trust-restoration gate.

\subsection{Screen attribution from the budget and scintillation constraints}
\label{sec:disc-screen-attribution}

% TODO(disc-screen-attribution): Explain how a screen attribution is made after
% the relevant Results entries exist. The logic should be:
% - first ask whether the Galactic contribution can explain the scattering;
% - then ask whether a validated foreground system can explain the scattering;
% - then use revalidated scintillation products to decide whether a one-screen
%   or two-screen interpretation is favored;
% - only then assign the remaining contribution to the host or to an unknown
%   intervening system.
%
% This subsection should turn the budget and scintillation products into a
% physical statement, not introduce new numbers. It should also name any
% sightline whose attribution remains coverage-limited or model-limited.
%
% The formal machinery for the one- vs two-screen decision is in
% Appendix~\ref{app:twoscreen}: the two-screen ACF product term and the
% m_tot -> sqrt(3) unresolved-regime signature (Eq. eq:sqrt3), the resolution
% power RP (Eq. eq:rp), and the screen-distance upper limit (Eq. eq:distance).
% When writing this subsection, cite Appendix~\ref{app:twoscreen} for the
% derivation. Reduced but detected MW modulation is evidence for partial
% resolution, not permission to saturate the inequality or quote a point
% distance. Any later distance posterior requires a calibrated RP likelihood
% with realization scatter and the MW-screen-distance uncertainty propagated.

One boundary condition for this attribution is now fixed by the CHIME-band
scintillation campaign (Section~\ref{sec:results-scintillation}). The
high-resolution FRB~20240203A product yields four qualified CHIME-band
decorrelation bandwidths, while one DSA-band point passes the full
qualification, $\gamma=0.446^{+0.239}_{-0.250}\,$MHz at $1328.24\,$MHz for
FRB~20220506D. Because those certified measurements belong to different
bursts, they do not supply a two-band frequency-scaling index or a common
$600\,$MHz $\Delta\nu_d$ anchor for any single sightline. The other 23 CHIME
products and the other DSA panels remain fit diagnostics rather than certified
bandwidth measurements.

The two-screen consistency test is not yet evaluable
(Table~\ref{tab:twoscreen-provisional}). The $\tau\,\Delta\nu_d$ product must
pair each DSA-band decorrelation bandwidth with a scattering time refit under
the same fixed convention---a single-exponential pulse-broadening function
with $\alpha$ fixed to 4---rather than with the free-$\alpha$ joint-fit
values, which serve the morphology audit and would mix two different $\tau$
definitions in one product. Those fixed-index refits do not yet exist. We
therefore report the seven otherwise eligible sightlines as pending rather
than compute screen products from an inconsistent $\tau$ track. Once the
refits exist, the calculation will evaluate $\tau(\nu)\,\Delta\nu_d$
separately at every clean narrow DSA component frequency, without imposing a
bandwidth power law, and compare it with the conservative common-screen
interval $0.1$--$2$. The largely uncalibrated DSA widths and the unavailable
full posterior covariance remain additional prerequisites. No two-screen,
screen-distance, or Galactic/host attribution is made from the current
products.

% Generated by scripts/build_provisional_propagation_tables.py; do not hand-edit.
\begin{deluxetable*}{lccccl}
\tablecaption{Provisional component-level screen-consistency products.\label{tab:twoscreen-provisional}}
\tablehead{\colhead{FRB} & \colhead{$N$} & \colhead{Median $\tau\Delta\nu_d$} & \colhead{Range} & \colhead{min$(P-\sigma_P)$} & \colhead{Inference}}
\startdata
20220310F & 1 & 3e+02 & 3e+02--3e+02 & 2.4e+02 & two-screen favored \\ 
20220506D & 4 & 4.2 & 1.5--26 & 0.45 & indeterminate \\ 
20221113A & 1 & 9.5 & 9.5--9.5 & 3.4 & two-screen favored \\ 
20230307A & 3 & 1.1e+02 & 36--2.5e+02 & 17 & two-screen favored \\ 
20230325A & 1 & 3.8e+02 & 3.8e+02--3.8e+02 & 2.4e+02 & two-screen favored \\ 
20230814B & 3 & 42 & 14--7.2e+02 & 9.8 & two-screen favored \\ 
20240122A & 3 & 99 & 46--4.7e+02 & 6.5 & two-screen favored \\ 
\enddata
\tablecomments{All values are best-so-far and provisional unless explicitly certified.}
\end{deluxetable*}


We next cross-match the provisional fitted-scattering values against the
validated foreground census of Section~\ref{sec:obs-fg}
(Table~\ref{tab:foreground-propagation-provisional}). This comparison does not
support a sample-wide association test: only seven propagation fits are
accepted, only two of the four sightlines with budget-eligible foreground
systems have accepted fits, and most foreground-negative sightlines lie outside
the deep-imaging footprints. Those zeros are censored rather than clean
foreground-free controls.

FRB~20230307A offers the strongest case for a partial foreground
contribution. Its sightline lies within the deep-imaging footprints and
intersects the confirmed cluster J115120.4+714435 at $b/R_{500}=0.83$ along
with three budget-eligible galaxy halos. The fiducial foreground model
predicts $\tau_{\rm int}=0.015$ ms, approximately 22\% of the provisional
$\tau_{1\,\mathrm{GHz}}=0.0676$ ms: the identified structures could supply
part of the observed broadening, but not all of it under the current model.
If a future two-screen result computed under the fixed-index convention
supports multiple propagation layers, this foreground complex would be a
candidate scattering screen; the present data neither assign a scintillation
screen to it nor determine a screen distance.

The confirmed halo approximately 101 kpc from the FRB~20220310F sightline
intersects it geometrically, but its fiducial scattering contribution is only
about 0.5\% of the provisional fitted value, disfavoring a dominant causal
role.
Conversely, FRBs~20221113A and 20230325A have close projected but inconclusive
candidates at 16 and 60 kpc, respectively; redshift validation is required
before either can enter the foreground budget. FRBs~20240229A and 20240203A
contain confirmed foreground halos but lack accepted propagation fits. For
FRB~20240229A the fiducial foreground prediction exceeds the retained
morphology-only fit value by a factor of about 11, which is a fit--budget tension
to resolve rather than evidence for causation.

\input{foreground_propagation_provisional_table}

\subsection{Event-by-event interpretation ledger}
\label{sec:disc-sightline-ledger}

Each co-detection should receive a short interpretation only after its relevant
budget, foreground, scintillation, fit, and energy entries have been restored.

% TODO(disc-zach / FRB 20220207C): Interpret after revalidation. Check whether
% the restored budget gives a foreground, host, Galactic, or unresolved
% attribution; state any coverage or fit-quality caveat.

% TODO(disc-whitney / FRB 20220310F): Interpret after revalidation. This slot
% should connect the multi-component modeling decision to any budget, screen, or
% energy statement that survives the new analysis.

% TODO(disc-oran / FRB 20220506D): Interpret after revalidation. Make the survey
% coverage status explicit before saying whether the lack of a foreground term
% is informative.

% TODO(disc-isha / FRB 20221113A): Interpret after revalidation. Separate the
% foreground-census meaning from any scattering-fit or energy result.

% TODO(disc-wilhelm / FRB 20221203A): Interpret after revalidation. Do not reuse
% the old scintillation or negative-residual language unless both strands are
% re-established.

% TODO(disc-phineas / FRB 20230307A): Interpret after revalidation. If the
% restored analysis finds multiple foreground systems or a cluster contribution,
% explain which ingredient dominates and which is secondary.

% TODO(disc-freya / FRB 20230325A): Interpret after revalidation. Keep any
% redshift-placeholder limitation explicit before using this event in
% distance-dependent statements.

% TODO(disc-johndoeii / FRB 20230814B): Interpret after revalidation. The
% legacy galaxy-interior classification did NOT survive the V4 census (no
% registered candidates; DM_int=0, note u — see FLITS PR #183); host z is the
% provisional internal 0.5535 (Verdi et al. in prep.), and the source is a
% repeater (earlier DSA-only detection, data unsaved);
% interpret from association, DM, and coverage information alone.

% TODO(disc-hamilton / FRB 20230913A): Interpret after revalidation. Treat the
% intervening-screen claim as open until the foreground, budget, scattering, and
% scintillation strands all independently pass trust restoration.

% TODO(disc-mahi / FRB 20240122A): Interpret after revalidation. State the
% coverage limitation before using any zero-foreground or host statement.

% TODO(disc-chromatica / FRB 20240203A): Interpret after revalidation. If a
% quality gate still fails, explain what can be said from association, DM, and
% foreground information alone.

% TODO(disc-casey / FRB 20240229A): Interpret after revalidation. State whether
% this sightline is used as a clean comparison case, a budget constraint, or an
% energy-only point.

\subsection{Burst morphology as a systematic on the turbulence spectrum}
\label{sec:disc-alpha}

% TODO(disc-alpha): Rebuild this subsection only after Section~\ref{sec:results-alpha}
% contains trusted fit results. The interpretation should distinguish:
% - what is learned about burst morphology;
% - what is learned about the scattering geometry;
% - what is learned about the turbulence spectrum;
% - which claims are method diagnostics rather than sky measurements.
%
% Do not quote old beta values, alpha limits, rail classes, class fractions, or
% the retired multiplicity-bias demonstration as results. If a revalidated
% multi-component case remains, explain why the extra component changes the
% physical interpretation rather than merely improving the fit statistic.

\subsection{Energy interpretation}
\label{sec:disc-energies}

% Scoped 2026-07-15 (V3 session; validation-v3-energetics-2026-07-15.md).
% Structure below is settled; sentences marked [FILL] take numbers from the
% regenerated (data-driven) burst_energies.json after owner V3 clearance.
% Do not quote legacy c0/gamma-derived energies or any per-band spectral
% index from the revoked fits anywhere in this subsection.

% Draft lead-in — uncomment with the five scoped paragraphs after V3 clears.
%The band-restricted energies of Section~\ref{sec:burst-energies} are
%deliberately conservative objects: each is the energy radiated within the
%two observed spectral windows, computed from the calibrated channel data with
%no spectral model, no extrapolation across the unobserved
%$0.8$--$1.3\,$GHz gap, and no fitted parameter inherited from the scattering
%analysis. Their limitations separate cleanly into five categories that we
%keep distinct.
%
% (1) Calibration. The absolute scale is catalog-anchored for three
% sightlines (~2x confidence) and radiometer-model-trusted for the rest;
% the 0.20-0.25 dex systematics dominate every row's error budget. No
% ranking of two rows closer than ~a factor of 2 is meaningful.
%
% (2) Spectral shape. In-band indices are unconstrained by design (narrow
% fractional bands; the DSA-band steepness the fits kept railing on is a
% model-family statement, not a measurement — cite the descriptive finding
% only if/when the C1 campaign re-establishes it). The two-band partition
% E_C vs E_D is the one shape-like quantity we CAN discuss: [FILL] state
% the observed E_C/E_D range; interpret order-unity scatter as evidence of
% per-burst spectral-envelope diversity across the octave gap, while noting
% that scintillation modulation and beam-model error enter that ratio and
% cannot be separated from intrinsic narrowbandedness until the two-band
% scintillation campaign concludes (Section~\ref{sec:disc-screen-attribution}).
%
% (3) Host redshift. Two provisional-z rows (FRBs 20230913A, 20240203A):
% any population-facing sentence must survive their removal; check and say
% so explicitly ([FILL] after table lands).
%
% (4) Selection. Co-detection requires threshold crossings at BOTH 600 MHz
% and 1.4 GHz: the sample is biased toward bursts that are simultaneously
% bright across the octave, i.e. against strongly banded events. The energy
% rows are therefore not a fair draw from either instrument's fluence
% function, and no luminosity-function normalization should be attempted
% from them. What CAN be said: [FILL] the span of ~3 decades in E_iso within
% a co-detected sample shows the joint-detection criterion still admits a
% wide intrinsic energy range, with the low end (~1e38-1e39 erg at
% z < 0.08 — zach, chromatica) illustrating that the local volume
% contributes the faint tail, consistent with the strong E-z selection
% correlation expected for flux-limited surveys. A qualitative comparison to
% the CHIME/FRB first-catalog energy distribution (Shin et al. 2023) is
% appropriate ONLY as an order-of-magnitude placement, with the band
% non-comparability caveat stated in the same sentence.
%
% (5) Forward link (lane E2 seed). Energies are the propagation-independent
% axis of the synthesis: fluence is conserved under scattering to first
% order, so E_iso vs (DM_host residual, tau, foreground class) scatter plots
% separate intrinsic brightness from path effects. With 8 rows this stays
% qualitative — frame as "no evident correlation" / "consistent with
% independence" unless something dramatic appears; a real statement needs
% the C1-re-fit tau values (gate: post-C, lane E2).

\subsection{Population limits and selection effects}
\label{sec:disc-population-limits}

% TODO(disc-population-limits): State which sample denominator applies to each
% claim. The twelve co-detections are the full association sample, but budget,
% foreground, scintillation, turbulence, and energy claims each have their own
% validated subset and exclusion reasons.
%
% Use this subsection to prevent over-reading the sample. Discuss redshift
% incompleteness, foreground-catalog coverage, fit-quality gates, and band
% selection effects only after the relevant Results text has named the affected
% sightlines.

\subsection{Implications for the next analysis pass}
\label{sec:disc-next-pass}

% TODO(disc-next-pass): Close the Discussion by naming the analysis that the
% revalidated results require next. Candidate topics include geometry selection,
% two-screen treatment, foreground-census sensitivity, burst-morphology
% modeling, or energy-selection rules. The closing paragraph should follow from
% the restored Results, not from the revoked campaign narrative.
